\documentclass[11pt,letterpaper]{article}
\usepackage[T1]{fontenc}
\usepackage{lmodern}
\usepackage[margin=1.02in]{geometry}
\usepackage{amsmath,amssymb,amsthm,booktabs,microtype,xurl}
\usepackage[hidelinks,pdftitle={Conditional ZF-independence of a 192-state Turing machine},
  pdfauthor={James Chang}]{hyperref}

\newtheorem{theorem}{Theorem}
\newtheorem{lemma}{Lemma}
\newcommand{\ZF}{\mathrm{ZF}}
\newcommand{\PA}{\mathrm{PA}}
\newcommand{\N}{\mathbb{N}}
\newcommand{\Con}{\operatorname{Con}}
\newcommand{\pair}[2]{\langle #1,#2\rangle}
\newcommand{\file}[1]{\nolinkurl{#1}}
\newcommand{\Left}{\mathsf{L}}
\newcommand{\Right}{\mathsf{R}}
\newcommand{\Halt}{H_{192}}
\newcommand{\phys}{\mathsf{Halt}_{192}}
\DeclareMathOperator{\AX}{AX}
\DeclareMathOperator{\MP}{MP}
\DeclareMathOperator{\SHIFT}{SHIFT}
\setlength{\emergencystretch}{2em}
\raggedbottom

\title{Conditional ZF-independence of a 192-state Turing machine}
\author{James Chang}
\date{}

\begin{document}
\maketitle

\begin{abstract}
We give an explicit binary Turing machine with 192 working states.
If first-order Zermelo--Fraenkel set theory is consistent, the machine
runs forever on a blank tape and the negation of its halting sentence
is unprovable in ZF. If ZF is arithmetically $\Sigma_1$-sound, the
halting sentence is also unprovable. The machine executes a
nine-register search for refutations in a coded calculus. A scanner
decodes instructions from the tape, and seven states perform the
register arithmetic. The simulation and the unprovability results
are formalized in Lean~4 for the transition table given in the paper.
\end{abstract}

\section{Introduction}

The running-time Busy Beaver function assigns to each state count
the longest blank-tape run among halting machines of that size. The
search for small machines with independent halting statements connects
this function with the strength of formal theories~\cite{yedidia-aaronson}.

Write $\ZF$ for first-order Zermelo--Fraenkel set theory and $\PA$
for first-order Peano arithmetic. We say that ZF is
\emph{arithmetically $\Sigma_1$-sound} if every closed $\Sigma_1$
arithmetic sentence whose natural-number interpretation is provable
in ZF is true in $\N$.

\begin{theorem}\label{thm:main}
There exists a closed sentence $\Halt$ in the language of set theory
whose interpretation in each model of ZF asserts that the machine of
Appendix~\ref{app:table} has a blank-tape halting history.
\begin{enumerate}
\item If ZF is consistent, the machine runs forever on the blank tape
and $\ZF\nvdash\neg\Halt$.
\item If ZF is arithmetically $\Sigma_1$-sound, then
$\ZF\nvdash\Halt$ and $\ZF\nvdash\neg\Halt$.
\end{enumerate}
\end{theorem}

The proof uses the register search in two ways. Over $\N$, a halting
computation gives a refutation in the calculus. A model of ZF
rules out such a refutation. Inside PA, a coded refutation of an
auxiliary arithmetic theory yields a halting history. The latter
implication allows a ZF proof of nonhalting to be turned into a
proof of the auxiliary theory's consistency in that theory itself.

Sections~\ref{sec:machine} and~\ref{sec:search} describe the machine
and the register search. Sections~\ref{sec:layout}--\ref{sec:simulation}
establish the simulation, and Section~\ref{sec:independence} proves
the nonhalting and unprovability results.

\subsection{Related work}

ZF207 and ZF209 give analogous conditional results over ZF for
their respective machines~\cite{zf207,zf209}. The two machines use
different decoders for a common binary program tree and share their
register arithmetic. The ZF192 table came from a separate search
(Section~\ref{sec:machine}). A neighboring bit distinguishes two
uses of a decoder state, and the seven-state arithmetic routine
shares a return scan between increment and positive decrement.
All three Lean proofs use arithmetic calibration and check their
own transition tables. For ZF192, the proof aligns the two control
graphs after unconditional jumps and uses induction in models of PA
for registers of arbitrary length.

\section{The machine and its halting sentence}\label{sec:machine}

\subsection{Transition rules}

The tape is indexed by $\mathbb{Z}$ and has alphabet $\{0,1\}$.
Initially all cells contain zero, the head is at coordinate $0$,
and the machine is in state $0$. The working states are
$0,\ldots,191$, and $192$ is the halting state.

A fresh search using Boolean satisfiability (SAT),
starting from Glazer's 230-state K230 machine~\cite{k230},
produced the ZF192 table. The search was independent of the ZF207 and ZF209
transition tables.

A row $(q,b,w,D,q')$ instructs the machine in state $q$ reading bit
$b$ to write $w$, move one cell in direction $D\in\{\Left,\Right\}$,
and enter state $q'$. Appendix~\ref{app:table} gives one row for
each working state and read bit. The only transitions to $192$
are $(0,1,1,\Right,192)$ and $(165,0,0,\Left,192)$. The first
cannot occur on the blank run because the initial read bit is $0$
and no transition enters state $0$. The 384 rows are also given in
\file{data/zf192.tm}, in order of state and then read bit.

\subsection{Coding configurations and histories}

We use the Cantor pairing function
\[
  \pair{x}{y}=\frac{(x+y)(x+y+1)}{2}+y,
\]
a bijection from $\N^2$ to $\N$. A configuration is represented by
$(q,\ell,r)\in\N^3$, with code $\pair{q}{\pair{\ell}{r}}$.
The binary digits of $\ell$ give the cells to the left of the head,
beginning with the nearest cell at the least significant bit.
Those of $r$ begin at the scanned cell and continue to the right.
Each configuration reached in finitely many steps has only finitely
many nonzero cells, so both codes are natural numbers.

A row reading $b$ applies when $r=b+2t$. A right move replaces
$(\ell,r)$ by $(w+2\ell,t)$. For a left move, write $\ell=c+2u$
with $c\in\{0,1\}$. The new pair is $(u,c+2w+4t)$. These
identities and the transition table define the arithmetic step relation.

A history lists the configuration codes in reverse chronological order.
The empty list has code $0$, and prepending a code $a$ to a list
with code $s$ gives $1+\pair{a}{s}$. A history of the blank run is
nonempty, starts chronologically at $(0,0,0)$, and satisfies the step
relation between consecutive configurations. It is a halting history
if its last configuration has no successor, which for a history of
this machine is equivalent to ending in state $192$. Internally, the
history check iterates the tail operation on a code $s$ for $s$ steps,
leaving the empty list fixed once reached. Each nonzero tail strictly
decreases, so this iteration checks the entire history in the model's
own natural numbers.

Let $\phys$ be the arithmetic assertion that such a halting history
exists. The formalization gives a $\Sigma_1$ presentation, with an
existentially quantified code for the computation. The sentence $\Halt$ is the
natural-number interpretation of the arithmetic halting assertion
in ZF. The formalization proves its equivalence to the corresponding
history predicate in each model, using that model's natural numbers.

\section{A register program for refutation search}\label{sec:search}

\subsection{Derivation rules}

Let $C$ be the calculus in Appendix~\ref{app:calculus}.
Its derivations use axiom introduction, modus ponens, and universal
generalization. Code $0$ denotes a tautology. Code $1$ denotes
$x_0\in x_0$, which is false under every assignment in a model of
ZF by the axiom of foundation. A \emph{$C$-refutation} is a finite
$C$-derivation containing code $1$.

\subsection{Enumeration of command words}

The search considers candidates $n=2,3,\ldots$ in order. It decodes
$n$ as a command and a pointer $k$ to the rest of the word. The
first component of the Cantor pair determines the command.
\begin{align*}
  \pair{0}{k} &\longmapsto (\SHIFT,k),\\
  \pair{1}{\pair{x}{k}} &\longmapsto (\MP(x),k),\\
  \pair{a+2}{\pair{y}{\pair{x}{k}}} &\longmapsto (\AX(a,x,y),k).
\end{align*}
A pointer $0$ ends the word. If a pointer is positive, the next
command is decoded from the pointer minus one. Every component
returned by unpairing is bounded by the input code, so the pointers
strictly decrease until they reach zero. Conversely, these rules
encode any finite nonempty command word.

\subsection{Execution of command words}

The store $(t,w,d,h)$ starts at $(0,0,0,0)$ and is carried from
one candidate to the next. Here $t$ and $h$ are formula codes,
$w$ codes a stack of formulas, and $d$ is an offset. Pushing $x$
replaces $w$ by $\pair{x}{w}$. Popping unpairs $w$, so a pop
from the empty stack returns code $0$ and leaves the stack empty.

\begin{description}
\item[$\SHIFT$] Halt if $t+h=1$. Otherwise move $h$ into $t$ and pop the
stack into $h$.
\item[$\MP(x)$] Set $v=d+x$ and test whether $h$ codes the implication
from $t$ to $v$. If it does, put $v$ in $t$ and push $0$ on the
stack. Otherwise put $0$ in $t$. In either case set $d=v$ and $h=0$.
\item[$\AX(a,x,y)$] Compute the seven axiom heads at $(a,d+x,h+y)$,
using $d$ as the fourth input to the seventh head. Push the heads,
the universal generalization of $h$ with binder $d$, and the incoming
value of $t$ in the order given in Appendix~\ref{app:calculus}.
If $a$ differs from both $d+x$ and $h+y$, put the corresponding pointed
hull axiom in $t$. Otherwise put $0$ in $t$. Set $d=h=0$.
\end{description}

The command rules are implemented by a program with nine unbounded
registers and 4,341 instruction sites. The instructions are increment,
test and decrement, unconditional jump, and halt. A test-and-decrement
instruction takes its zero branch without changing the register. On
the positive branch it subtracts one. The program decoded from the
tape has 3,379 live sites. The correspondence is established in
Section~\ref{sec:simulation}.

\subsection{Correctness of the search}

\begin{lemma}\label{lem:search}
Over $\N$, the register search halts from its initial store if and
only if there is a $C$-refutation.
\end{lemma}

\begin{proof}
Initially $t$, $h$, and all stack entries are derivable in $C$.
The command rules preserve this invariant. In particular, $\MP$
introduces a nonzero formula only by modus ponens, and $\AX$
adds axiom instances and a universal generalization. The halting
test $t+h=1$ requires one of $t$ and $h$ to be $1$, so a halt
implies a $C$-refutation.

Conversely, fix a $C$-derivation of $1$. Its compiled word must
work with the store left by the preceding candidate. To prepare
that store, the word begins with $\AX(0,0,0)$, eight applications
of $\SHIFT$, and $\MP(0)$. This sequence acts by
\[
  (t,w,d,h)\longmapsto (0,\pair{t}{w},0,0).
\]
The first command places seven implication codes and a universal
code above the incoming $t$ on the stack. The shifts bring the
incoming $t$ to the top and leave the universal code in $h$.
The tags of the generated codes preclude halting during these
shifts. The final $\MP(0)$ compares this universal code with an
implication code, so it clears the formula registers without
changing the stack.

Appendix~\ref{app:compiler} gives a command sequence for each
inference rule. Inductively, a sequence for a formula code $b$,
started from any store $(t,w,0,0)$, either halts or ends at
$(b,w,0,0)$. For a derivation of $1$, put this sequence after
the preparation and append $\SHIFT$. The final test halts
if the sequence has returned $1$. The word starts with $\AX$,
so its candidate number is at least $2$. Each earlier candidate
either halts or finishes and advances the counter. The search
therefore halts on this word or on an earlier candidate.
\end{proof}

\section{Encoding the register program on the tape}\label{sec:layout}

\subsection{Register configurations}

A \emph{source configuration} on the tape records the current
instruction address and the nine register values. The machine is
in state $157$ at the address root. In the relative coordinates
of this configuration the head is at $0$. The tape contents are given
in these coordinates.
\begin{center}
\begin{tabular}{@{}ll@{}}
\toprule
Coordinates & Tape contents \\
\midrule
$0,-1,\ldots,-G+1$ & Complemented address bits followed by ones \\
$-G,-G-1$ & Two zeros forming the gate, with $G\geq39$ \\
$-G-2,-G-3,\ldots$ & Nine register blocks, followed by zeros \\
$1$ through $8$ & A zero, six ones, then a zero \\
Greater than $8$ & Arbitrary bits of finite support \\
\bottomrule
\end{tabular}
\end{center}
Starting just to the left of the gate and reading leftward, a
register with value $v$ is represented by $1^{v+1}0$. The initial
one distinguishes a zero register from the two consecutive zeros
at the gate.

\subsection{The initial computation}

The first transition writes a one and moves left into state $1$.
After $44{,}181$ steps the machine is in state $157$, with the head
at coordinate $-121$ in the original tape coordinates. The gate
zeros are at $-207$ and $-208$, giving $G=86$. The address is
$100000$, which designates site $22$, and the register values are
$(2,0,0,0,0,0,0,0,0)$. The tape has $186$ ones, including those
beyond the zero cell at relative coordinate $8$. Lean verifies the
entire tape at this checkpoint by replaying the initial computation
in five intervals on packed tape words.

\section{Implementation of the register instructions}\label{sec:macros}

The decoder chooses among ten increment selectors and ten decrement
selectors. A selector moves left through the register file, counting
zero separators until it reaches the chosen block. The arithmetic
routines use states $79,96,111,118,124,125,134$. They return toward
the address in state $117$ after an increment or a positive decrement,
and in state $129$ after testing a zero register.

\subsection{The return scan}\label{sec:marker}

Both successful operations move the gate one cell to the left. The
shared return scan runs toward the address. Throughout the scan,
a zero immediately to the left of the head marks the boundary
between cells already shifted and those still to be copied.
Starting in state $125$, the machine takes one of the following
paths according to the bits at and to the right of the head.

\begin{center}
\begin{tabular}{@{}lll@{}}
\toprule
Bits at the head & States along the path & Effect \\
\midrule
$1$ & $125,125,134,96,111,125$ & Move a body bit left \\
$01$ & $125,134,96,111,125$ & Move a separator and a body bit left \\
$00$ & $125,134,117$ & Complete the shift at the gate \\
\bottomrule
\end{tabular}
\end{center}

The first path writes the body bit into the marker cell and
makes the cell it vacated the new marker. The second moves the
pair $01$ one cell left and advances the marker past the pair.
The pattern $00$ marks the gate. The last path completes the
displacement by turning one of its original cells into padding.
Each of the three paths is checked directly against the transition
table, with arbitrary tape on either side of the cells used.

An induction on the number of body cells proves that the scan
crosses a register block. Repeating the argument for the finitely
many intervening blocks gives the return to the gate. The first
induction uses an arithmetic formula, so it also applies to
register lengths in nonstandard models of PA.

\subsection{Test and decrement}

From the first one of the selected block, state $118$ moves left
into state $111$. If the head reads the separator, the register
is zero. State $111$ leaves that separator intact and enters
state $124$, which scans right. The scan crosses an isolated
separator through states $118$ and $111$, then resumes in $124$.
At the gate, state $118$ moves into state $129$. This branch
leaves the register file and gate intact.

If state $111$ reads a second one, it replaces that one by zero
and returns to the first one in state $125$. The new zero is the
marker for the return scan of Section~\ref{sec:marker}.
The scan removes one body cell from the chosen block, leaves the
other register values intact, and moves the gate one cell left.
The return begins in state $117$.

\subsection{Increment}

The increment first scans left through the selected block and
the blocks farther from the address. State $79$ writes a one and
continues left, entering state $96$ when it reads a zero. If state
$96$ reads a one, it writes zero and resumes the scan in state $79$.
When state $96$ reaches the zero beyond the last separator,
it writes one and turns right into state $111$. Two cells suffice
to set up the marker and enter state $125$.

Together, the outward scan and the return insert one body cell in
the selected block. They preserve the other eight values and the
zero tape beyond the file. The gate moves one cell left and the
return begins in state $117$.

\section{Simulation of the register program}\label{sec:simulation}

\subsection{Local decoding and return paths}

The decoder sometimes uses a neighboring bit to distinguish two
situations in the same state. For instance, when state $128$ reads
zero, it moves left into state $136$. If the new cell is zero,
states $136$ and $137$ bring the head back through the original
cell and one position to its right. If the new cell is one,
state $136$ moves one cell further left into state $143$. Both
paths preserve the tape.

Every decoder path is checked in the window from $-38$ through $8$
in the relative coordinates of Section~\ref{sec:layout}. At one of
the 3,379 live instructions, the path must halt, reach the root for
the next address, or enter the appropriate arithmetic selector at
the left boundary. The checks also cover 4,339 arithmetic returns,
including both branches of each test-and-decrement instruction.
They verify the resulting state, head position, and window contents.
Each path reads only cells in the window, so the same transitions
hold for any tape outside it.

The selector reaches the gate by moving left across the remaining
ones without changing them. After arithmetic, state $117$ or $129$
crosses the padding to the right in the same way. Joining the two
scans to the finite paths and arithmetic routines gives a positive
number of tape steps for every contextual instruction. The gate
distance remains at least $39$.

For example, site $22$, reached by the initial computation, tests
register $8$. The zero branch leads to site $24$ and the positive
branch to site $23$. This register initially contains zero.
The zero branch takes $248$ steps, returns the head to $-121$ in
state $157$, and changes only the cell at $-125$.

\subsection{Correspondence of control graphs}

The accompanying source gives both finite graphs. The register
program is specified by a binary tree of primitive instructions
in \file{TreeData.lean}. Its leaves are numbered from left to right,
and the path to a leaf determines the targets of its fall-through
and control-flow edges. In \file{data/source.json}, the entry at
index $i$ of \file{sites} gives the tape address of contextual site
$i$, and the entry in \file{cfg} gives its instruction and successors.
The table has 3,689 slots, of which 310 are unused.

Each graph has 3,037 sites for arithmetic or halting. The correspondence pairs these sites in increasing order
within each program. For a jump site, follow its unconditional
successors to the first arithmetic or halt site and use that
site's partner. Such a jump prefix has length at most three.

If the register-program values are $(r_0,\ldots,r_8)$, the tape
stores them in the order
\[
  (r_6,r_2,r_0,r_4,r_3,r_7,r_8,r_5,r_1).
\]
Lean checks that paired arithmetic instructions operate on the
corresponding register
and that their successor sites correspond after following any
unconditional jumps. The check also identifies the paired halt
sites and the initial configurations.

\subsection{Standard and internal simulation}

\begin{lemma}\label{lem:simulation}
Over $\N$, the tape machine halts on the blank tape if and only
if the register search halts from its initial configuration.
In every model of PA, a coded $C$-refutation yields a coded
blank-tape halting history.
\end{lemma}

\begin{proof}
An unconditional jump in the register program leaves the
corresponding tape configuration unchanged. An arithmetic
instruction takes a positive number of tape steps to execute the
matching operation and the subsequent jumps in the decoded
program. The halt instruction leads to state $192$. Thus a finite
register computation ending in halt yields a halting computation
of the tape machine. An infinite register computation contains
infinitely many arithmetic instructions, because the jump prefixes
are bounded. It therefore gives arbitrarily long tape computations
from the initial configuration. Determinism rules out a halting
computation from that configuration.

The internal argument uses the arithmetic definition of the
relation between a register configuration and its tape
representation. In a model of PA, the compiler turns a coded
$C$-refutation into a coded register computation ending in the
halt instruction. Induction on the coded computation gives a
halting history for the tape machine. The induction applies to
histories of nonstandard length.
\end{proof}

\section{Nonhalting and independence over ZF}\label{sec:independence}

\subsection{Soundness of the coded calculus}

The axioms of $C$ hold under their set-theoretic interpretation
in every model of ZF. The logical heads represent valid inference
templates. The set-existence axiom expresses the pointed hull
principle of Appendix~\ref{app:calculus}. To construct a hull for
$R$ in ZF, start with $\{v_0\}$ if $R(u_0,v_0)$ holds and with
the empty set otherwise. At each stage, consider the subsets
$u$ of members of the current set for which some $R(u,v)$ holds.
For each such $u$ adjoin all $R$-witnesses of least rank. These
form a set, and replacement and union collect them over the set
of requests. The union of the stages indexed by $\omega$
contains the required witnesses and every member has an
$R$-predecessor. The inference rules of $C$ preserve truth under
every assignment.

If ZF is consistent, it has a model by first-order completeness.
Code $1$ is false under every assignment in that model, so no
$C$-refutation exists over $\N$. Lemmas~\ref{lem:search}
and~\ref{lem:simulation} now imply that the blank-tape run is infinite.

\subsection{Auxiliary theories}

We use an arithmetic theory $T$ whose refutations can be translated
to $C$ and whose theorems include the arithmetic consequences of ZF.
First define a set-language theory $U$ from derivations in $C$.
Let $\operatorname{Cert}_C(n,k)$ mean that $n$ codes a $C$-derivation
containing formula code $k$, and let $\chi(\theta)$ be the
translation into $C$ of a closed set-language sentence $\theta$.

For an arithmetic sentence $\sigma$, write $J_e(\sigma)$ for its
interpretation in the natural numbers defined relative to the marker
$e$. Write $K(\sigma)$ for the assertion that this interpretation
holds at every empty marker.
\[
  K(\sigma)=\forall e\bigl(\operatorname{Empty}(e)\to J_e(\sigma)\bigr),
  \qquad
  \operatorname{Empty}(e)\equiv\forall x\,(x\notin e).
\]
Define $\top^{(0)}=\top$ and $\top^{(n+1)}=\top\land\top^{(n)}$.
The two axiom sets consist of the following sentences.
\[
\begin{array}{ll}
 U & \theta\land\top^{(n)}
     \quad\text{when }\operatorname{Cert}_C(n,\chi(\theta)),\\[2pt]
 T & \sigma\land\top^{(n)}
     \quad\text{when }\operatorname{Cert}_C(n,\chi(K(\sigma))).
\end{array}
\]
The repeated conjunctions record the certificate index $n$ in each
axiom. To decide whether a sentence is an axiom, extract this index
and check the derivation it codes. The two axiom predicates are
$\Delta_1$ over $I\Sigma_1$, the fragment of PA with induction for
$\Sigma_1$ formulas. In particular, $I\Sigma_1$ proves the
equivalence of the $\Sigma_1$ and $\Pi_1$ descriptions of each
predicate.

\subsection{Interpretation and transfer of provability}

The interpretation first identifies objects $x,y$ with the same
membership behavior in both arguments of $\in$. The relation on
the quotient is again written $\in$. In a model of $U$, the
pointed hull principle holds on this quotient. Applying it to
the empty relation gives an empty object, since every member of
a hull must have a predecessor under that relation.

Fix an empty object $e$ in the quotient. An object is marked if
it contains $e$. For marked objects $x,y$, write $x\sim_e y$ when
they are equal or are related by a set-coded bisimulation for
marked membership. A bisimulation matches each marked member
of $x$ with an equal or related marked member of $y$, and
conversely. The second quotient is by $\sim_e$. In this quotient,
$[x]\in[y]$ means that $y$ has a marked member equivalent to $x$.
The ZF interpretation restricts to the elements of transitive
sets that are internally well-founded in this second quotient.
The pointed hull construction supplies its set-existence
axioms.

For the interpretation $J_e$, take the natural numbers of this
ZF model. The formal arithmetic translation uses the projection
that sends a representative of a natural number to that number
and all other objects to zero. Quantifying over representatives
therefore gives quantification over the interpreted natural
numbers.

In a model of ZF, soundness of $C$ gives $K(\sigma)$ for every
certified axiom of $T$. Evaluating these assertions at an empty
marker gives a model of $T$. Thus consistency of ZF implies
consistency of $T$.

Let $\sigma^\omega$ denote the usual interpretation of an
arithmetic sentence $\sigma$ in the natural numbers of set theory.
We use the following transfer of provability.
\begin{equation}\label{eq:transfer}
  \ZF\vdash\sigma^\omega \quad\Longrightarrow\quad T\vdash\sigma .
\end{equation}
Suppose $\ZF\vdash\sigma^\omega$. In every model of $U$, each
empty marker induces the marked ZF interpretation, so
$K(\sigma)$ holds in that model. By completeness,
$U\vdash K(\sigma)$. The translation of this derivation into $C$
uses the stored certificate for each axiom of $U$ and introduces its
conjunctions with $\top$. The result is a code $n$ such that
$\operatorname{Cert}_C(n,\chi(K(\sigma)))$. Hence
$\sigma\land\top^{(n)}$ is an axiom of $T$, yielding $T\vdash\sigma$.
In particular, $T$ extends PA because ZF proves the natural-number
interpretations of the PA axioms.

\subsection{Unprovability of halting and nonhalting}

In a model of PA, a coded $T$-refutation gives a coded $U$-refutation
by the marked interpretation and the existence of an empty marker.
The stored certificates at the axiom uses give a coded
$C$-refutation. By Lemma~\ref{lem:simulation} this yields a halting
history for the tape machine. Arithmetic completeness gives the
corresponding implication in PA. Its contrapositive, interpreted in
the natural numbers of ZF, gives the following, where
$\Con(T)$ is the consistency sentence for the axiom predicate of $T$.
\begin{equation}\label{eq:consistency}
  \ZF\vdash\neg\Halt\longrightarrow\Con(T)^\omega .
\end{equation}

If ZF is consistent, then so is $T$. A ZF proof of $\neg\Halt$
would yield a ZF proof of $\Con(T)^\omega$ by
\eqref{eq:consistency}. Applying \eqref{eq:transfer} would give
$T\vdash\Con(T)$, contrary to G\"odel's second incompleteness
theorem~\cite{godel}. Thus $\ZF\nvdash\neg\Halt$.

If ZF is arithmetically $\Sigma_1$-sound, it is consistent and the
actual blank-tape run is infinite. A proof of $\Halt$ in ZF would
give a proof of the natural-number interpretation of the $\Sigma_1$
halting assertion. By soundness this would give a halting history
over $\N$. Hence $\ZF\nvdash\Halt$, completing the proof
of Theorem~\ref{thm:main}.

\section{Formalization in Lean}\label{sec:reproduction}

The standalone \texttt{ZF192.Challenge} defines the literal machine,
its history semantics, and both claims. The proved declarations in
\texttt{ZF192.Solution} are
\begin{quote}\small
\file{PalomarZF192.exact192Claim}\\
\file{PalomarZF192.nonhalting_of_consistent}.
\end{quote}
The first gives the sentence $\Halt$, its semantics in every ZF model,
and both conditional unprovability results. The second proves nonhalting
under consistency of ZF. Both use the exact table and depend only on
Lean's \file{propext}, \file{Classical.choice}, and \file{Quot.sound} axioms.

From the accompanying Lean project, run the following with Lean 4.35.0
to obtain the optional Mathlib cache, build the proof, and audit the axioms.
\begin{verbatim}
lake exe cache get
bash scripts/check.sh --full
\end{verbatim}
On Linux with Bubblewrap and user namespaces, the standalone
comparison and NanoDa replay can be run with
\begin{verbatim}
lake comparator --config comparator.json
\end{verbatim}

\begin{thebibliography}{9}
\raggedright
\bibitem{godel}
Kurt G\"odel.
\newblock \"Uber formal unentscheidbare S\"atze der \emph{Principia
  Mathematica} und verwandter Systeme I.
\newblock \emph{Monatshefte f\"ur Mathematik und Physik}, 38:173--198,
  1931. \href{https://doi.org/10.1007/BF01700692}{doi:10.1007/BF01700692}.

\bibitem{yedidia-aaronson}
Adam Yedidia and Scott Aaronson.
\newblock A relatively small Turing machine whose behavior is independent
  of set theory.
\newblock \href{https://arxiv.org/abs/1605.04343}{arXiv:1605.04343}, 2016.
\bibitem{zf207}
James Chang.
\newblock Nonhalting and ZF-independence for a 207-state Turing machine.
\newblock 2026.
\href{https://github.com/jamesyc/BusyBeaver/blob/9437ef874e1c76a0ec45330554e895157d56c604/ZF207/paper/zf207.pdf}{Paper}
and
\href{https://github.com/jamesyc/BusyBeaver/tree/9437ef874e1c76a0ec45330554e895157d56c604/ZF207}{Lean formalization}.

\bibitem{zf209}
James Chang.
\newblock Nonhalting and ZF-independence for a 209-state Turing machine.
\newblock 2026.
\href{https://github.com/jamesyc/BusyBeaver/blob/286324f9ad106bbeab32b70404e7a82be2cd660d/ZF209/paper/zf209.pdf}{Paper}
and
\href{https://github.com/jamesyc/BusyBeaver/tree/286324f9ad106bbeab32b70404e7a82be2cd660d/ZF209}{Lean formalization}.

\bibitem{k230}
Elliot Glazer.
\newblock K230: a 230-state Turing machine for the consistency of ZFC.
\newblock 2026.
\href{https://github.com/elliotglazer/k230-rocq/tree/3142e4dd980f0be12f4abb636c4ae2528afc17eb}{Machine and Rocq formalization}.

\end{thebibliography}

\appendix
\section{Numerical presentation of the calculus}\label{app:calculus}

\subsection{Formula codes and axiom heads}

Write $I(a,b)$ for the code of implication, $M(a,b)$ for membership,
$E(a,b)$ for equality, and $Q_x(a)$ for universal quantification with
binder $x$. The exact operations are
\[
 I(a,b)=4\pair{a}{b},\quad M(a,b)=4\pair{a}{b}+1,\quad
 Q_x(a)=4\pair{a}{x}+2,\quad E(a,b)=4\pair{a}{b}+3 .
\]
In the formulas below, $a\to b$ abbreviates $I(a,b)$ and associates to
the right. Put $F=1$. The seven heads at inputs $(a,d,h;u)$ are
\begin{align*}
 A_1 &= Q_d(a)\to Q_d(a\to h)\to Q_d(h),\\
 A_2 &= Q_d(E(h,d)\to F)\to F,\\
 A_3 &= E(h,d)\to M(a,h)\to M(a,d),\\
 A_4 &= M(h,d)\to E(h,a)\to M(a,d),\\
 A_5 &= Q_d(Q_a(h))\to Q_a(Q_d(h)),\\
 A_6 &= \bigl(Q_d(Q_d(h)\to F)\to F\bigr)\to h,\\
 A_7 &= \bigl((h\to u)\to d\bigr)\to(d\to h)\to a\to h .
\end{align*}
Head $A_7$ is an axiom when $u\le d$. The axiom set also includes $0$.

\subsection{The pointed hull axiom}

The axiom $G(a,d,h)$ is present when $a\ne d$ and $a\ne h$.
Its code is built from the same operations.
\begin{align*}
 B &= Q_h\bigl(a\to (Q_h(a\to M(h,a)\to F)\to F)\bigr),\\
 S &= Q_d\bigl(Q_a(M(a,d)\to M(a,h))\to B\bigr),\\
 P &= Q_h\bigl(M(h,a)\to (Q_d(a\to S\to F)\to F)\bigr),\\
 G(a,d,h)
   &= \bigl((E(h,d)\to Q_a(E(h,d)))\to
          (M(h,d)\to Q_a(M(h,d)))\to {}\\[-2pt]
   &\hspace{3.8em} Q_a\bigl(P\to((a\to M(h,a))\to F)\bigr)\bigr)\to F .
\end{align*}
All binders act on the displayed numerical labels. This is a coded
set-existence schema. Its interpretation, with parameters, supplies a
pointed hull for a definable relation $R$ and a pair $(u_0,v_0)$.
Such a hull is a set $C_0$ satisfying
\[
\begin{split}
 R(u_0,v_0)&\ \Longrightarrow\ v_0\in C_0,\\
 z\in C_0&\ \Longrightarrow\ \exists p\ R(p,z),\\
 z\in C_0,\quad u\subseteq z,\quad \exists v\ R(u,v)
    &\ \Longrightarrow\ \exists w\in C_0\ R(u,w).
\end{split}
\]

\subsection{The axiom command}

For $\AX(a,x,y)$, evaluate the seven heads at $(a,d+x,h+y;d)$.
The new stack begins, from the top, with
\[
  (A_1,A_2,A_3,A_4,A_5,A_6,A_7,Q_d(h),t),
\]
followed by the prior stack. The generalization $Q_d(h)$ uses the
incoming values of $d$ and $h$.

\subsection{Compilation of derivations}\label{app:compiler}

Write $S=\SHIFT$ and use juxtaposition for concatenation of
command words. The following words implement stack operations.
\begin{align*}
 P &= \AX(0,0,0) S^8\MP(0),\\
 L &= \AX(0,0,0) S^{10},\\
 D &= \MP(1)\AX(1,0,0)S^9\MP(0),\\
 J_i &= \begin{cases}
   \text{the empty word}, & i=0,\\
   \MP(i), & i>0.
 \end{cases}
\end{align*}
The word $P$ pushes $t$ and clears $t,d,h$. If
$w=\pair{x}{w'}$, the word $L$ sets $(w,d,h)=(w',0,x)$
while preserving $t$. From $(t,w,0,0)$ the word $D$ returns
$(0,w,0,0)$, and from $(0,w,0,0)$ the word $J_i$ sets $d=i$.

For an axiom, the following words leave its code in $t$, preserve
$w$, and return $d=h=0$ when started with $d=h=0$.
\[
\begin{array}{ll}
 0 & D,\\
 A_j(a,d,h)\ (1\le j\le 6) & D\,\AX(a,d,h)\,S^{j+1}\,L^{8-j},\\
 A_7(a,d,h;u)\ (u\le d) & D\,J_u\,\AX(a,d-u,h)\,S^8\,L,\\
 G(a,d,h)\ (a\ne d,\ a\ne h) & D\,\AX(a,d,h)\,L^9 .
\end{array}
\]
The guard on $A_7$ makes the increment $d-u$ well-defined, and
$J_u$ supplies the incoming offset of that command.

For generalization, first run the word for $a$. If it returns,
append $P\,J_i\,S\,\AX(i,0,0)\,S^9$ to obtain $Q_i(a)$.
For modus ponens, run the word for $I(a,b)$, then $P$, then
the word for $a$, then $P$ again. At this point the first two
stack entries are $a$ and $I(a,b)$. The suffix
\[
  S^2\,\MP(b)\,\AX(b,0,0)\,S^{10}
\]
either halts when a shift encounters code $1$, or returns
$(b,w,0,0)$ with the original stack $w$. These store
transformations give the induction used in Lemma~\ref{lem:search}.

\clearpage
\section{Complete transition table}\label{app:table}

The file \file{data/zf192.tm} has SHA-256 digest
\begin{quote}\small
\file{be1db915a6f40e37a601d3cb54bc0cf0d9597bdf0c75e97098cf4c201c6d5585}.
\end{quote}

Each entry is $(\text{write},\text{direction},\text{next state})$.
The direction is the head movement in the original tape coordinates.
The initial state is $0$ and the halting state is $192$.

\subsection*{States 0--95}

\begin{center}
\small
\setlength{\tabcolsep}{4pt}
\renewcommand{\arraystretch}{1.10}
\begin{tabular}{@{}rll@{\hspace{1.1em}}rll@{\hspace{1.1em}}rll@{}}
\toprule
State & Read 0 & Read 1 & State & Read 0 & Read 1 & State & Read 0 & Read 1 \\
\midrule
0 & $(1,\mathsf{L},1)$ & $(1,\mathsf{R},192)$ & 32 & $(0,\mathsf{R},11)$ & $(0,\mathsf{L},51)$ & 64 & $(0,\mathsf{L},84)$ & $(1,\mathsf{L},85)$ \\
1 & $(0,\mathsf{L},2)$ & $(1,\mathsf{L},3)$ & 33 & $(0,\mathsf{L},52)$ & $(1,\mathsf{L},33)$ & 65 & $(1,\mathsf{R},86)$ & $(1,\mathsf{L},87)$ \\
2 & $(0,\mathsf{L},3)$ & $(1,\mathsf{L},4)$ & 34 & $(0,\mathsf{L},53)$ & $(1,\mathsf{L},54)$ & 66 & $(0,\mathsf{L},50)$ & $(1,\mathsf{L},66)$ \\
3 & $(0,\mathsf{L},5)$ & $(1,\mathsf{L},6)$ & 35 & $(1,\mathsf{R},55)$ & $(0,\mathsf{L},26)$ & 67 & $(0,\mathsf{L},26)$ & $(1,\mathsf{L},67)$ \\
4 & $(0,\mathsf{L},7)$ & $(1,\mathsf{L},8)$ & 36 & $(0,\mathsf{L},51)$ & $(1,\mathsf{L},36)$ & 68 & $(0,\mathsf{L},88)$ & $(0,\mathsf{L},89)$ \\
5 & $(0,\mathsf{L},9)$ & $(1,\mathsf{L},10)$ & 37 & $(0,\mathsf{L},47)$ & $(1,\mathsf{L},34)$ & 69 & $(0,\mathsf{L},67)$ & $(1,\mathsf{L},69)$ \\
6 & $(0,\mathsf{L},11)$ & $(1,\mathsf{L},12)$ & 38 & $(1,\mathsf{L},26)$ & $(0,\mathsf{L},56)$ & 70 & $(0,\mathsf{L},56)$ & $(1,\mathsf{L},70)$ \\
7 & $(0,\mathsf{L},13)$ & $(1,\mathsf{L},14)$ & 39 & $(0,\mathsf{R},13)$ & $(1,\mathsf{L},40)$ & 71 & $(1,\mathsf{R},90)$ & $(1,\mathsf{L},72)$ \\
8 & $(0,\mathsf{L},15)$ & $(1,\mathsf{L},16)$ & 40 & $(0,\mathsf{L},18)$ & $(1,\mathsf{R},57)$ & 72 & $(0,\mathsf{L},86)$ & $(0,\mathsf{L},89)$ \\
9 & $(0,\mathsf{L},17)$ & $(1,\mathsf{L},18)$ & 41 & $(0,\mathsf{L},58)$ & $(1,\mathsf{L},12)$ & 73 & $(1,\mathsf{R},55)$ & $(0,\mathsf{L},89)$ \\
10 & $(0,\mathsf{L},19)$ & $(1,\mathsf{L},12)$ & 42 & $(0,\mathsf{L},59)$ & $(1,\mathsf{L},42)$ & 74 & $(0,\mathsf{L},91)$ & $(1,\mathsf{L},74)$ \\
11 & $(0,\mathsf{L},20)$ & $(1,\mathsf{L},21)$ & 43 & $(0,\mathsf{L},60)$ & $(1,\mathsf{L},61)$ & 75 & $(0,\mathsf{L},92)$ & $(1,\mathsf{L},6)$ \\
12 & $(0,\mathsf{L},22)$ & $(1,\mathsf{L},20)$ & 44 & $(0,\mathsf{L},62)$ & $(1,\mathsf{L},10)$ & 76 & $(0,\mathsf{L},93)$ & $(1,\mathsf{R},94)$ \\
13 & $(0,\mathsf{L},23)$ & $(1,\mathsf{L},24)$ & 45 & $(0,\mathsf{L},63)$ & $(1,\mathsf{L},64)$ & 77 & $(1,\mathsf{L},69)$ & $(0,\mathsf{L},50)$ \\
14 & $(0,\mathsf{L},25)$ & $(0,\mathsf{L},26)$ & 46 & $(0,\mathsf{L},65)$ & $(0,\mathsf{L},66)$ & 78 & $(0,\mathsf{R},95)$ & $(0,\mathsf{L},70)$ \\
15 & $(0,\mathsf{L},27)$ & $(1,\mathsf{L},6)$ & 47 & $(1,\mathsf{L},67)$ & $(0,\mathsf{L},66)$ & 79 & $(1,\mathsf{L},96)$ & $(1,\mathsf{L},79)$ \\
16 & $(0,\mathsf{L},20)$ & $(1,\mathsf{R},28)$ & 48 & $(0,\mathsf{L},68)$ & $(0,\mathsf{L},56)$ & 80 & $(1,\mathsf{R},17)$ & $(1,\mathsf{L},97)$ \\
17 & $(1,\mathsf{R},17)$ & $(0,\mathsf{L},26)$ & 49 & $(0,\mathsf{L},69)$ & $(1,\mathsf{L},49)$ & 81 & $(0,\mathsf{L},82)$ & $(1,\mathsf{R},98)$ \\
18 & $(0,\mathsf{L},29)$ & $(1,\mathsf{L},30)$ & 50 & $(0,\mathsf{L},36)$ & $(1,\mathsf{L},50)$ & 82 & $(0,\mathsf{L},99)$ & $(1,\mathsf{L},12)$ \\
19 & $(0,\mathsf{L},31)$ & $(1,\mathsf{L},32)$ & 51 & $(0,\mathsf{L},70)$ & $(1,\mathsf{L},51)$ & 83 & $(0,\mathsf{L},100)$ & $(1,\mathsf{L},101)$ \\
20 & $(1,\mathsf{R},33)$ & $(1,\mathsf{R},34)$ & 52 & $(0,\mathsf{L},49)$ & $(1,\mathsf{L},52)$ & 84 & $(1,\mathsf{L},102)$ & $(0,\mathsf{L},66)$ \\
21 & $(0,\mathsf{L},35)$ & $(0,\mathsf{L},36)$ & 53 & $(0,\mathsf{L},71)$ & $(1,\mathsf{R},72)$ & 85 & $(0,\mathsf{L},68)$ & $(0,\mathsf{L},103)$ \\
22 & $(0,\mathsf{L},37)$ & $(1,\mathsf{L},38)$ & 54 & $(0,\mathsf{L},73)$ & $(0,\mathsf{L},74)$ & 86 & $(1,\mathsf{L},26)$ & $(0,\mathsf{L},50)$ \\
23 & $(1,\mathsf{R},17)$ & $(1,\mathsf{L},39)$ & 55 & $(1,\mathsf{L},33)$ & $(0,\mathsf{L},66)$ & 87 & $(1,\mathsf{R},17)$ & $(1,\mathsf{L},17)$ \\
24 & $(0,\mathsf{L},3)$ & $(1,\mathsf{L},5)$ & 56 & $(0,\mathsf{L},74)$ & $(1,\mathsf{L},56)$ & 88 & $(1,\mathsf{R},90)$ & $(1,\mathsf{L},104)$ \\
25 & $(0,\mathsf{L},40)$ & $(1,\mathsf{L},41)$ & 57 & $(0,\mathsf{L},75)$ & $(1,\mathsf{L},76)$ & 89 & $(0,\mathsf{L},33)$ & $(1,\mathsf{L},89)$ \\
26 & $(0,\mathsf{L},42)$ & $(1,\mathsf{L},26)$ & 58 & $(0,\mathsf{L},77)$ & $(1,\mathsf{L},78)$ & 90 & $(1,\mathsf{R},105)$ & $(0,\mathsf{L},50)$ \\
27 & $(0,\mathsf{L},43)$ & $(1,\mathsf{L},44)$ & 59 & $(0,\mathsf{L},79)$ & $(1,\mathsf{L},59)$ & 91 & $(0,\mathsf{L},106)$ & $(1,\mathsf{L},91)$ \\
28 & $(0,\mathsf{L},18)$ & $(1,\mathsf{L},45)$ & 60 & $(0,\mathsf{L},80)$ & $(1,\mathsf{L},41)$ & 92 & $(0,\mathsf{L},9)$ & $(1,\mathsf{L},41)$ \\
29 & $(0,\mathsf{L},46)$ & $(1,\mathsf{L},34)$ & 61 & $(0,\mathsf{L},81)$ & $(1,\mathsf{L},82)$ & 93 & $(0,\mathsf{L},107)$ & $(1,\mathsf{L},28)$ \\
30 & $(0,\mathsf{L},47)$ & $(1,\mathsf{L},48)$ & 62 & $(0,\mathsf{L},83)$ & $(1,\mathsf{L},57)$ & 94 & $(0,\mathsf{L},108)$ & $(1,\mathsf{L},109)$ \\
31 & $(1,\mathsf{L},49)$ & $(0,\mathsf{L},50)$ & 63 & $(0,\mathsf{L},10)$ & $(1,\mathsf{L},31)$ & 95 & $(0,\mathsf{L},54)$ & $(1,\mathsf{R},110)$ \\
\bottomrule
\end{tabular}
\end{center}

\clearpage
\subsection*{States 96--191}

\begin{center}
\small
\setlength{\tabcolsep}{4pt}
\renewcommand{\arraystretch}{1.10}
\begin{tabular}{@{}rll@{\hspace{1.1em}}rll@{\hspace{1.1em}}rll@{}}
\toprule
State & Read 0 & Read 1 & State & Read 0 & Read 1 & State & Read 0 & Read 1 \\
\midrule
96 & $(1,\mathsf{R},111)$ & $(0,\mathsf{L},79)$ & 128 & $(0,\mathsf{L},136)$ & $(1,\mathsf{R},137)$ & 160 & $(1,\mathsf{R},55)$ & $(1,\mathsf{L},164)$ \\
97 & $(0,\mathsf{R},112)$ & $(1,\mathsf{R},24)$ & 129 & $(1,\mathsf{R},135)$ & $(1,\mathsf{R},129)$ & 161 & $(0,\mathsf{L},16)$ & $(1,\mathsf{L},165)$ \\
98 & $(0,\mathsf{L},113)$ & $(1,\mathsf{L},114)$ & 130 & $(0,\mathsf{L},138)$ & $(1,\mathsf{L},75)$ & 162 & $(1,\mathsf{R},127)$ & $(1,\mathsf{L},90)$ \\
99 & $(0,\mathsf{L},18)$ & $(0,\mathsf{L},26)$ & 131 & $(0,\mathsf{L},139)$ & $(1,\mathsf{L},25)$ & 163 & $(0,\mathsf{L},162)$ & $(0,\mathsf{L},36)$ \\
100 & $(0,\mathsf{L},23)$ & $(1,\mathsf{L},115)$ & 132 & $(0,\mathsf{L},140)$ & $(1,\mathsf{L},41)$ & 164 & $(0,\mathsf{L},166)$ & $(1,\mathsf{L},167)$ \\
101 & $(0,\mathsf{L},25)$ & $(1,\mathsf{L},100)$ & 133 & $(0,\mathsf{L},141)$ & $(1,\mathsf{L},142)$ & 165 & $(0,\mathsf{L},192)$ & $(1,\mathsf{L},156)$ \\
102 & $(0,\mathsf{L},89)$ & $(1,\mathsf{L},102)$ & 134 & $(1,\mathsf{R},117)$ & $(1,\mathsf{L},96)$ & 166 & $(0,\mathsf{L},110)$ & $(1,\mathsf{L},168)$ \\
103 & $(0,\mathsf{L},66)$ & $(1,\mathsf{L},103)$ & 135 & $(0,\mathsf{L},129)$ & $(0,\mathsf{R},128)$ & 167 & $(0,\mathsf{L},169)$ & $(1,\mathsf{L},170)$ \\
104 & $(0,\mathsf{L},55)$ & $(1,\mathsf{L},116)$ & 136 & $(0,\mathsf{R},137)$ & $(1,\mathsf{L},143)$ & 168 & $(0,\mathsf{L},171)$ & $(0,\mathsf{L},70)$ \\
105 & $(1,\mathsf{R},117)$ & $(1,\mathsf{R},105)$ & 137 & $(0,\mathsf{R},128)$ & $(1,\mathsf{R},128)$ & 169 & $(0,\mathsf{L},77)$ & $(1,\mathsf{L},64)$ \\
106 & $(0,\mathsf{L},118)$ & $(1,\mathsf{L},106)$ & 138 & $(0,\mathsf{L},144)$ & $(1,\mathsf{R},7)$ & 170 & $(0,\mathsf{L},172)$ & $(1,\mathsf{L},173)$ \\
107 & $(0,\mathsf{L},119)$ & $(1,\mathsf{L},120)$ & 139 & $(0,\mathsf{L},120)$ & $(1,\mathsf{L},39)$ & 171 & $(1,\mathsf{R},90)$ & $(1,\mathsf{L},174)$ \\
108 & $(0,\mathsf{L},97)$ & $(1,\mathsf{L},121)$ & 140 & $(0,\mathsf{L},39)$ & $(1,\mathsf{L},6)$ & 172 & $(0,\mathsf{L},175)$ & $(1,\mathsf{L},176)$ \\
109 & $(0,\mathsf{L},122)$ & $(1,\mathsf{R},123)$ & 141 & $(0,\mathsf{L},145)$ & $(1,\mathsf{L},146)$ & 173 & $(1,\mathsf{R},127)$ & $(0,\mathsf{L},103)$ \\
110 & $(0,\mathsf{L},11)$ & $(1,\mathsf{L},20)$ & 142 & $(0,\mathsf{L},147)$ & $(1,\mathsf{L},126)$ & 174 & $(0,\mathsf{L},177)$ & $(1,\mathsf{L},178)$ \\
111 & $(0,\mathsf{R},124)$ & $(0,\mathsf{R},125)$ & 143 & $(0,\mathsf{R},117)$ & $(1,\mathsf{L},148)$ & 175 & $(0,\mathsf{L},90)$ & $(0,\mathsf{L},42)$ \\
112 & $(0,\mathsf{L},92)$ & $(1,\mathsf{L},99)$ & 144 & $(0,\mathsf{L},149)$ & $(1,\mathsf{L},147)$ & 176 & $(0,\mathsf{L},179)$ & $(0,\mathsf{L},91)$ \\
113 & $(0,\mathsf{L},108)$ & $(1,\mathsf{L},5)$ & 145 & $(0,\mathsf{L},150)$ & $(1,\mathsf{L},151)$ & 177 & $(0,\mathsf{L},180)$ & $(1,\mathsf{L},28)$ \\
114 & $(0,\mathsf{L},126)$ & $(1,\mathsf{L},18)$ & 146 & $(0,\mathsf{L},78)$ & $(1,\mathsf{L},81)$ & 178 & $(0,\mathsf{L},181)$ & $(0,\mathsf{L},70)$ \\
115 & $(0,\mathsf{L},60)$ & $(1,\mathsf{R},112)$ & 147 & $(0,\mathsf{L},18)$ & $(1,\mathsf{L},92)$ & 179 & $(0,\mathsf{R},32)$ & $(0,\mathsf{L},102)$ \\
116 & $(1,\mathsf{R},127)$ & $(0,\mathsf{L},50)$ & 148 & $(0,\mathsf{R},117)$ & $(1,\mathsf{L},152)$ & 180 & $(0,\mathsf{L},182)$ & $(1,\mathsf{L},20)$ \\
117 & $(0,\mathsf{L},128)$ & $(1,\mathsf{R},117)$ & 149 & $(0,\mathsf{L},80)$ & $(1,\mathsf{L},6)$ & 181 & $(0,\mathsf{L},183)$ & $(1,\mathsf{L},1)$ \\
118 & $(0,\mathsf{R},129)$ & $(1,\mathsf{L},111)$ & 150 & $(0,\mathsf{L},153)$ & $(1,\mathsf{L},154)$ & 182 & $(0,\mathsf{L},10)$ & $(1,\mathsf{L},156)$ \\
119 & $(0,\mathsf{L},130)$ & $(1,\mathsf{L},131)$ & 151 & $(0,\mathsf{L},155)$ & $(1,\mathsf{L},156)$ & 183 & $(0,\mathsf{L},123)$ & $(1,\mathsf{L},184)$ \\
120 & $(0,\mathsf{L},40)$ & $(1,\mathsf{R},27)$ & 152 & $(0,\mathsf{L},157)$ & $(1,\mathsf{L},143)$ & 184 & $(0,\mathsf{L},98)$ & $(1,\mathsf{L},185)$ \\
121 & $(0,\mathsf{L},18)$ & $(1,\mathsf{L},82)$ & 153 & $(1,\mathsf{R},127)$ & $(1,\mathsf{L},151)$ & 185 & $(0,\mathsf{L},186)$ & $(1,\mathsf{L},187)$ \\
122 & $(0,\mathsf{L},108)$ & $(1,\mathsf{L},132)$ & 154 & $(0,\mathsf{L},32)$ & $(1,\mathsf{L},55)$ & 186 & $(0,\mathsf{L},94)$ & $(1,\mathsf{L},188)$ \\
123 & $(0,\mathsf{L},133)$ & $(1,\mathsf{R},109)$ & 155 & $(1,\mathsf{R},117)$ & $(1,\mathsf{R},95)$ & 187 & $(0,\mathsf{L},189)$ & $(1,\mathsf{L},138)$ \\
124 & $(0,\mathsf{R},118)$ & $(1,\mathsf{R},124)$ & 156 & $(0,\mathsf{L},158)$ & $(1,\mathsf{L},159)$ & 188 & $(0,\mathsf{L},112)$ & $(1,\mathsf{L},25)$ \\
125 & $(0,\mathsf{R},134)$ & $(1,\mathsf{L},125)$ & 157 & $(0,\mathsf{L},160)$ & $(1,\mathsf{L},161)$ & 189 & $(0,\mathsf{L},190)$ & $(1,\mathsf{L},149)$ \\
126 & $(0,\mathsf{L},121)$ & $(1,\mathsf{R},98)$ & 158 & $(0,\mathsf{L},162)$ & $(1,\mathsf{L},155)$ & 190 & $(0,\mathsf{L},191)$ & $(1,\mathsf{L},147)$ \\
127 & $(1,\mathsf{R},127)$ & $(1,\mathsf{R},135)$ & 159 & $(1,\mathsf{R},90)$ & $(1,\mathsf{L},163)$ & 191 & $(0,\mathsf{L},115)$ & $(1,\mathsf{L},140)$ \\
\bottomrule
\end{tabular}
\end{center}

\end{document}
